%======================================================================
% c) Comparacion con un programa comercial
%======================================================================
\newpage
\section*{c) Comparaci\'on con un programa comercial de elementos finitos}

%----------------------------------------------------------------------
\begin{desarrollo}{1. Programa y tipo de an\'alisis}

Como programa comercial us\'e \textbf{Autodesk Inventor Nastran} (solver
Autodesk Nastran), con un an\'alisis \emph{est\'atico lineal}. A diferencia de
mi c\'odigo, el programa parte de la pieza CAD (\texttt{parte1.ipt}), genera la
malla autom\'aticamente y entrega los campos ya post-procesados.

Hay dos formas razonables de modelar la placa en un programa comercial:
\begin{itemize}[itemsep=3pt]
  \item \textbf{Elementos 2D (placa/l\'amina) en la superficie media}, con
  espesor $t$. Con carga solo en el plano, una l\'amina trabaja como membrana y
  el resultado es equivalente a tensi\'on plana. Es lo m\'as comparable con mi
  c\'odigo. Los elementos t\'ipicos de Nastran son \texttt{CQUAD8} (cuadril\'atero
  de 8 nodos, \emph{serend\'ipito}, el Q8 de mi c\'odigo) y \texttt{CTRIA6}
  (tri\'angulo de 6 nodos, el T6).
  \item \textbf{Elementos s\'olidos 3D} (tetraedros de 10 nodos) en todo el
  volumen de la pieza. No necesita la hip\'otesis de tensi\'on plana, pero tiene
  muchos m\'as GDL y la tensi\'on puede variar un poco a trav\'es del espesor.
\end{itemize}

\textbf{Ojo con los GDL de las l\'aminas.} Los elementos de l\'amina de Nastran
tienen 6 GDL por nodo (3 traslaciones y 3 rotaciones), no 2 como en elasticidad
plana. Con carga en el plano los GDL fuera del plano no trabajan, pero hay que
restringirlos (o al menos uno por cuerpo) para que la rigidez no quede singular.
\end{desarrollo}

%----------------------------------------------------------------------
\begin{desarrollo}{2. Modelo en el programa comercial}

Para que la comparaci\'on sea justa, el modelo comercial debe resolver
\emph{el mismo problema} que el apartado a): misma geometr\'ia, material,
espesor, condiciones de borde y carga. As\'i las diferencias que queden se
deben solo a la malla y al tipo de elemento.

En el programa se model\'o un cuarto de la placa, igual que en el apartado a),
con simetr\'ia en los bordes $x=0$ e $y=0$ y la tracci\'on $S=100$ MPa aplicada
en el borde $x=L$ (figura de condiciones de borde). La malla es autom\'atica y
se refin\'o en el borde importante, el del agujero, donde se usaron
\textbf{100 elementos}, porque ah\'i est\'a el gradiente de tensiones y el
m\'aximo.

\begin{center}
\figuraT{T3/figuras/nastran_malla.png}{0.48}{Malla del programa comercial, con detalle del refinamiento junto al agujero.}\hfill
\figuraT{T3/figuras/nastran_cb.png}{0.48}{Condiciones de borde y carga en el programa comercial.}
\end{center}
\end{desarrollo}

%----------------------------------------------------------------------
\begin{desarrollo}{3. Resultados del programa comercial}

\begin{center}
\figuraT{T3/figuras/nastran_sxx.png}{0.48}{$\sigma_{xx}$ (componente normal en la direcci\'on de la carga).}\hfill
\figuraT{T3/figuras/nastran_vonmises.png}{0.48}{Tensi\'on de von Mises.}

\vspace{6pt}
\figuraT{T3/figuras/nastran_desplazamiento.png}{0.48}{Desplazamiento en $x$.}
\end{center}

El m\'aximo de $\sigma_{xx}$ es $319.3$ MPa y aparece en el borde del agujero,
sobre el eje perpendicular a la carga (punto B); el de von Mises es
$313.9$ MPa en el mismo lugar, y el desplazamiento m\'aximo en $x$ es
$0.102$ mm, en el borde cargado.
\end{desarrollo}

%----------------------------------------------------------------------
\begin{desarrollo}{4. Comparaci\'on: Kirsch, c\'odigo creado y programa comercial}

\begin{center}
\renewcommand{\arraystretch}{1.25}
\begin{tabular}{@{}lccc@{}}
\toprule
Variable & Kirsch (placa infinita) & C\'odigo creado (malla 2) & Inventor Nastran\\
\midrule
$\sigma_{xx}(B)$ [MPa] & $300.0$ & $304.0$ & $319.3$\\
$K_t=\sigma_{xx}(B)/S$ & $3.00$ & $3.04$ & $3.19$\\
$u_x$ m\'aximo [mm] & $0.10066$ & $0.10127$ & $0.102$\\
$\sigma_{VM,\max}$ [MPa] & $300.0$ & $301.0$ & $313.9$\\
\bottomrule
\end{tabular}
\end{center}

\textbf{Qu\'e esperar.} Si los dos modelos (mi c\'odigo y Nastran) resuelven el mismo problema
(misma placa finita, mismas condiciones de borde), deber\'ian parecerse
\emph{entre s\'i} m\'as que a Kirsch, porque los dos tienen el efecto de ancho y
largo finitos discutido en b). Las diferencias entre mi c\'odigo y Nastran vienen
de:
\begin{itemize}[itemsep=3pt]
  \item \textbf{Malla}: la m\'ia es estructurada y graduada a mano; la del
  programa es autom\'atica, no estructurada y en general con muchos m\'as
  elementos.
  \item \textbf{Post-proceso de tensiones}: yo eval\'uo $\bm\sigma$ en los nodos
  y promedio entre elementos; el programa puede mostrar valores en las esquinas,
  en los centroides o suavizados, y el m\'aximo en el borde del agujero depende de
  eso.
\end{itemize}

\textbf{Discusi\'on.}
\begin{itemize}[itemsep=3pt]
  \item Los tres m\'etodos coinciden en lo esencial: concentraci\'on de
  tensiones en B, con $K_t=3.00$ (Kirsch), $3.04$ (c\'odigo creado) y $3.19$
  (Nastran). El c\'odigo creado queda a $1.3$\,\% de Kirsch y Nastran a $6.3$\,\%.
  \item Nastran difiere de mi c\'odigo en $5$\,\% en $K_t$ y en $4$\,\% en
  von Mises, m\'as que mi c\'odigo respecto de Kirsch. No es lo que esperaba
  (que los dos modelos se parecieran m\'as entre s\'i). Lo atribuyo a la malla
  autom\'atica y no estructurada, y a que la tensi\'on m\'axima en el borde del
  agujero es un valor puntual muy sensible al tama\~no local del elemento y a
  c\'omo el programa post-procesa las tensiones, aun con los $100$ elementos en
  ese borde.
  \item Los desplazamientos, en cambio, casi no dependen de la malla:
  $0.102$ mm en Nastran contra $0.1013$ mm en mi c\'odigo ($0.7$\,\% de
  diferencia), porque son una cantidad global.
  \item \textbf{Programa comercial:} importa la geometr\'ia CAD, mallado
  autom\'atico y post-proceso listo, lo que ahorra mucho tiempo. \textbf{C\'odigo
  propio:} se controla todo el modelo, se sabe exactamente qu\'e se integra y
  c\'omo se calculan las tensiones, y se puede refinar de forma sistem\'atica
  para estudiar la convergencia, como en b).
\end{itemize}
\end{desarrollo}
